\chapter{Release 1 : Socle citoyen}
\label{chap:release1}

% Affichage robuste : le rapport reste compilable avant l'export PNG depuis Draw.io.
\providecommand{\releasefigure}[4]{%
\begin{figure}[H]
    \centering
    \IfFileExists{#1}{%
        \includegraphics[width=#2\textwidth]{#1}%
    }{%
        \fbox{\parbox[c][4.2cm][c]{0.88\textwidth}{\centering
        \textbf{Diagramme à exporter depuis Draw.io}\\[0.3cm]
        \texttt{\detokenize{#4}}\\
        vers\\
        \texttt{\detokenize{#1}}}}
    }%
    \caption{#3}
\end{figure}}

\section{Introduction de la Release}
\label{sec:r1_introduction}

La Release 1 constitue le socle citoyen d'EcoRewind. Conformément à la planification établie au chapitre~\ref{chap:analyse_specification_besoins}, elle regroupe trois sprints complémentaires. Le premier sécurise l'accès à la plateforme et la gestion des profils. Le deuxième met à disposition la carte, les points de collecte et les consignes locales de tri. Le troisième introduit le fil communautaire, les interactions, les témoignages et les notifications.

Cette organisation fournit un incrément utilisable : le citoyen peut créer son compte, accéder aux informations de tri de son territoire et participer à la communauté. Les fonctionnalités éducatives, les récompenses et l'intégration IoT sont réservées aux releases suivantes.

\section{Sprint 1 : Authentification, profils et gestion des rôles}
\label{sec:r1_sprint1}

\subsection{Objectifs et backlog du Sprint 1}

L'objectif du Sprint 1 est de mettre en place une gestion sécurisée des identités. Le backend fournit l'inscription, la connexion par JWT, le renouvellement de session, la récupération du mot de passe, l'authentification sociale et la vérification multifacteur. Le modèle \texttt{User} centralise le rôle, le code QR unique, l'avatar et les paramètres MFA.

\begin{longtable}{|>{\centering\arraybackslash}p{1cm}|p{9.2cm}|>{\centering\arraybackslash}p{1.8cm}|>{\centering\arraybackslash}p{1.8cm}|}
\caption{Backlog du Sprint 1 -- Accès, profils et rôles}\label{tab:r1_s1_backlog}\\
\hline
\rowcolor{gray!35}\textbf{ID} & \textbf{User Story} & \textbf{Priorité} & \textbf{État}\\\hline
\endfirsthead
\hline\rowcolor{gray!35}\textbf{ID} & \textbf{User Story} & \textbf{Priorité} & \textbf{État}\\\hline
\endhead
1.1 & En tant que visiteur, je veux créer et vérifier mon compte afin d'accéder aux services personnalisés. & Élevée & Réalisée\\\hline
1.2 & En tant qu'utilisateur, je veux me connecter, renouveler ma session et me déconnecter de manière sécurisée. & Élevée & Réalisée\\\hline
1.3 & En tant qu'utilisateur, je veux consulter et modifier mon profil, mon avatar et mon QR personnel. & Élevée & Réalisée\\\hline
1.4 & En tant qu'utilisateur, je veux récupérer mon mot de passe et activer la MFA afin de renforcer la sécurité de mon compte. & Élevée & Réalisée\\\hline
1.5 & En tant qu'administrateur, je veux gérer les comptes et leurs rôles selon mes autorisations. & Élevée & Réalisée\\\hline
\end{longtable}

\subsection{Raffinement des cas d'utilisation}

La figure~\ref{fig:r1_s1_uc} place les acteurs à l'extérieur de la frontière d'EcoRewind. L'authentification est une précondition des opérations privées et non un sous-cas inclus artificiellement. La vérification MFA étend la connexion uniquement lorsque cette option est activée. La vérification du compte fait partie du parcours d'inscription.

\releasefigure{image/usecase_sprint1_access.png}{0.95}{Diagramme de cas d'utilisation du Sprint 1}{usecase_sprint1_access.drawio}
\label{fig:r1_s1_uc}

\subsection{Description textuelle du cas « S'authentifier »}

\begin{longtable}{|p{3.2cm}|p{10.2cm}|}
\caption{Description du cas d'utilisation « S'authentifier »}\label{tab:r1_uc_auth}\\
\hline\rowcolor{gray!20}\textbf{Élément} & \textbf{Description}\\\hline
Acteur principal & Utilisateur disposant d'un compte EcoRewind.\\\hline
Préconditions & Le compte existe et il est actif.\\\hline
Postconditions & Un jeton d'accès et un jeton de renouvellement sont produits, sauf si une vérification MFA reste requise.\\\hline
Scénario nominal & 1. L'utilisateur saisit son courriel et son mot de passe. \newline 2. Flutter transmet les données à \texttt{POST /token}. \newline 3. Le routeur recherche l'utilisateur et vérifie le mot de passe haché. \newline 4. Le backend synchronise les données utiles avec Firebase. \newline 5. Il retourne les jetons, le rôle, l'identifiant et le QR personnel. \newline 6. Flutter ouvre l'espace correspondant au rôle.\\\hline
Alternatives & A1 : identifiants invalides, réponse HTTP 401. \newline A2 : MFA active, retour d'un \texttt{mfa\_token}, saisie du code TOTP puis appel de \texttt{/auth/mfa/verify}.\\\hline
\end{longtable}

\subsection{Diagrammes de séquence}

Le diagramme de la figure~\ref{fig:r1_s1_seq_auth} respecte le flux réel de \texttt{POST /token}. Le fragment \texttt{alt} distingue l'échec, la connexion directe et la branche MFA. Les retours sont représentés en pointillés et les appels synchrones en traits pleins.

\releasefigure{image/sequence_sprint1_authentication.png}{0.98}{Diagramme de séquence « S'authentifier »}{sequence_sprint1_authentication.drawio}
\label{fig:r1_s1_seq_auth}

La création du compte vérifie l'unicité du courriel, hache le mot de passe avec bcrypt, persiste l'utilisateur et génère le QR par la valeur par défaut du modèle. La synchronisation Firebase est non bloquante.

\releasefigure{image/sequence_sprint1_signup.png}{0.98}{Diagramme de séquence « Créer un compte »}{sequence_sprint1_signup.drawio}
\label{fig:r1_s1_seq_signup}

\subsection{Diagramme de classes du Sprint 1}

La figure~\ref{fig:r1_s1_class} ne crée pas de sous-classes artificielles pour les rôles : ceux-ci sont stockés dans l'attribut \texttt{role} de \texttt{User}. \texttt{OTPCode} représente les codes temporaires utilisés pour l'inscription ou la réinitialisation.

\releasefigure{image/class_sprint1_access.png}{0.90}{Diagramme de classes du Sprint 1}{class_sprint1_access.drawio}
\label{fig:r1_s1_class}

\subsection{Scénarios de validation}

Les scénarios de recette portent sur l'inscription avec courriel unique, le refus d'identifiants invalides, la production des JWT, la branche MFA, la protection d'une ressource privée et le contrôle des rôles. Ils constituent des critères de validation ; leur présence dans ce tableau ne remplace pas l'exécution des tests automatisés.

\section{Sprint 2 : Carte, points de collecte et consignes de tri}
\label{sec:r1_sprint2}

\subsection{Objectifs et backlog du Sprint 2}

Ce sprint rend l'information de tri directement exploitable. Le citoyen consulte les points sur une carte OpenStreetMap, recherche un centre et filtre les résultats. L'administrateur maintient les points de collecte, tandis que l'intercommunalité gère les consignes locales de son territoire.

\begin{longtable}{|>{\centering\arraybackslash}p{1cm}|p{9.2cm}|>{\centering\arraybackslash}p{1.8cm}|>{\centering\arraybackslash}p{1.8cm}|}
\caption{Backlog du Sprint 2 -- Points de collecte et consignes}\label{tab:r1_s2_backlog}\\
\hline\rowcolor{gray!35}\textbf{ID} & \textbf{User Story} & \textbf{Priorité} & \textbf{État}\\\hline
\endfirsthead
\hline\rowcolor{gray!35}\textbf{ID} & \textbf{User Story} & \textbf{Priorité} & \textbf{État}\\\hline
\endhead
2.1 & En tant que visiteur ou citoyen, je veux consulter les points de collecte sur une carte. & Élevée & Réalisée\\\hline
2.2 & En tant que citoyen, je veux rechercher et filtrer les points puis consulter leurs détails. & Élevée & Réalisée\\\hline
2.3 & En tant qu'administrateur, je veux créer, modifier et supprimer un point de collecte. & Élevée & Réalisée\\\hline
2.4 & En tant qu'intercommunalité, je veux gérer les consignes locales par territoire, ville et type de déchet. & Élevée & Réalisée\\\hline
\end{longtable}

\subsection{Raffinement des cas d'utilisation}

La consultation générale des points est distinguée de la vue centralisée protégée de l'intercommunalité. « Filtrer les points » et « Afficher les détails » étendent la consultation car ces comportements sont déclenchés seulement à la demande. Les opérations d'administration sont associées directement aux acteurs autorisés.

\releasefigure{image/usecase_sprint2_collection_points.png}{0.96}{Diagramme de cas d'utilisation du Sprint 2}{usecase_sprint2_collection_points.drawio}
\label{fig:r1_s2_uc}

\subsection{Description textuelle du cas « Consulter les points de collecte »}

\begin{longtable}{|p{3.2cm}|p{10.2cm}|}
\caption{Description du cas d'utilisation « Consulter les points de collecte »}\label{tab:r1_uc_points}\\
\hline\rowcolor{gray!20}\textbf{Élément} & \textbf{Description}\\\hline
Acteurs & Visiteur et citoyen.\\\hline
Préconditions & L'application peut joindre l'API ; aucune authentification n'est imposée pour la consultation publique.\\\hline
Postconditions & Les marqueurs correspondant aux données disponibles sont affichés sur la carte.\\\hline
Scénario nominal & 1. L'acteur ouvre la carte. \newline 2. Flutter demande la liste des points. \newline 3. Le routeur interroge \texttt{CollectionPoint}. \newline 4. L'API retourne les coordonnées, types, adresse, horaires, statut et niveau de remplissage. \newline 5. Flutter place les marqueurs.\\\hline
Alternatives & A1 : liste vide, affichage d'un état vide. \newline A2 : erreur réseau, affichage d'une erreur avec possibilité de relancer. \newline A3 : filtres renseignés, mise à jour des résultats visibles.\\\hline
\end{longtable}

\subsection{Diagramme de classes du Sprint 2}

\texttt{CollectionPoint} et \texttt{LocalInstruction} sont deux concepts indépendants : aucune association persistante n'est inventée entre eux. La seule association représentée relie l'auteur à ses consignes grâce à la clé étrangère \texttt{created\_by}.

\releasefigure{image/class_sprint2_collection_points.png}{0.94}{Diagramme de classes du Sprint 2}{class_sprint2_collection_points.drawio}
\label{fig:r1_s2_class}

\subsection{Diagramme de séquence}

La consultation publique ne contient pas de fragment \texttt{ref S'authentifier}. Le fragment \texttt{alt} décrit les trois résultats possibles : liste disponible, liste vide ou erreur technique.

\releasefigure{image/sequence_sprint2_collection_points.png}{0.98}{Diagramme de séquence « Consulter les points de collecte »}{sequence_sprint2_collection_points.drawio}
\label{fig:r1_s2_seq}

\section{Sprint 3 : Fil communautaire, témoignages et notifications}
\label{sec:r1_sprint3}

\subsection{Objectifs et backlog du Sprint 3}

Le Sprint 3 introduit la participation citoyenne. Une publication est enregistrée immédiatement avec le statut \texttt{pending\_review}, puis analysée par une tâche de fond. Un contenu sûr peut être publié automatiquement ; un contenu signalé reste soumis à la décision humaine. L'IA ne prononce donc pas seule un rejet définitif dans le flux principal.

\begin{longtable}{|>{\centering\arraybackslash}p{1cm}|p{9.2cm}|>{\centering\arraybackslash}p{1.8cm}|>{\centering\arraybackslash}p{1.8cm}|}
\caption{Backlog du Sprint 3 -- Communauté et notifications}\label{tab:r1_s3_backlog}\\
\hline\rowcolor{gray!35}\textbf{ID} & \textbf{User Story} & \textbf{Priorité} & \textbf{État}\\\hline
\endfirsthead
\hline\rowcolor{gray!35}\textbf{ID} & \textbf{User Story} & \textbf{Priorité} & \textbf{État}\\\hline
\endhead
3.1 & En tant que citoyen, je veux consulter et créer des publications avec un média éventuel. & Élevée & Réalisée\\\hline
3.2 & En tant que citoyen, je veux aimer, commenter, répondre et enregistrer une publication. & Élevée & Réalisée\\\hline
3.3 & En tant que citoyen, je veux soumettre un témoignage ou proposer un point de collecte. & Moyenne & Réalisée\\\hline
3.4 & En tant qu'utilisateur, je veux recevoir et marquer comme lues mes notifications. & Élevée & Réalisée\\\hline
3.5 & En tant qu'administrateur, je veux examiner les contenus signalés et décider de leur publication ou rejet. & Élevée & Réalisée\\\hline
\end{longtable}

\subsection{Raffinement des cas d'utilisation}

L'ajout d'un média étend la création d'une publication, car il reste facultatif. Aimer, commenter et enregistrer sont des spécialisations comportementales optionnelles de l'interaction avec une publication. La modération automatique est incluse dans le traitement d'une nouvelle publication ; la décision humaine intervient seulement pour un contenu signalé.

\releasefigure{image/usecase_sprint3_community.png}{0.97}{Diagramme de cas d'utilisation du Sprint 3}{usecase_sprint3_community.drawio}
\label{fig:r1_s3_uc}

\subsection{Description textuelle du cas « Publier une action citoyenne »}

\begin{longtable}{|p{3.2cm}|p{10.2cm}|}
\caption{Description du cas d'utilisation « Publier une action citoyenne »}\label{tab:r1_uc_post}\\
\hline\rowcolor{gray!20}\textbf{Élément} & \textbf{Description}\\\hline
Acteur & Citoyen authentifié.\\\hline
Préconditions & Le jeton d'accès est valide ; la description et les informations attendues respectent le schéma \texttt{PostCreate}.\\\hline
Postconditions & La publication est persistée en \texttt{pending\_review} et une tâche de modération est planifiée.\\\hline
Scénario nominal & 1. Le citoyen rédige son contenu et ajoute éventuellement un média. \newline 2. Flutter téléverse le fichier si nécessaire et obtient son URL. \newline 3. Flutter appelle \texttt{POST /posts}. \newline 4. L'API persiste la publication. \newline 5. Elle répond immédiatement puis lance \texttt{\_run\_ai\_moderation} en arrière-plan. \newline 6. Le statut et les informations d'audit sont mis à jour. \newline 7. Une notification informe l'auteur.\\\hline
Alternatives & A1 : score sûr, statut \texttt{published}. \newline A2 : contenu signalé ou erreur du pipeline, maintien en \texttt{pending\_review} pour arbitrage humain. \newline A3 : données invalides ou média non téléversé, la requête est rejetée avant création.\\\hline
\end{longtable}

\subsection{Diagramme de classes du Sprint 3}

\texttt{Like} et \texttt{SavedPost} sont représentées comme des classes d'association entre \texttt{User} et \texttt{Post}. La relation réflexive de \texttt{Comment} modélise les réponses. Les liens de notification ne sont dessinés comme associations que lorsqu'ils correspondent à de véritables clés étrangères.

\releasefigure{image/class_sprint3_community.png}{0.99}{Diagramme de classes du Sprint 3}{class_sprint3_community.drawio}
\label{fig:r1_s3_class}

\subsection{Diagramme de séquence}

Le fragment \texttt{opt [média fourni]} précède la création. La réponse HTTP est renvoyée avant la modération asynchrone. Le fragment \texttt{alt} distingue la publication automatique et la révision humaine ; il ne présente pas un rejet automatique comme décision finale du pipeline principal.

\releasefigure{image/sequence_sprint3_publish_post.png}{0.99}{Diagramme de séquence « Publier une action citoyenne »}{sequence_sprint3_publish_post.drawio}
\label{fig:r1_s3_seq}

\section{Implémentation de la Release 1}
\label{sec:r1_implementation}

L'application cliente est développée avec Flutter. Elle communique avec les routeurs FastAPI par HTTP et conserve les éléments de session nécessaires. Les données métier sont persistées par SQLAlchemy et les évolutions du schéma sont gérées par Alembic. Firebase complète le backend pour les usages configurés, notamment la synchronisation et les notifications mobiles.

La Release 1 mobilise principalement les modules \texttt{auth}, \texttt{users}, \texttt{collection\_points}, \texttt{intercommunality}, \texttt{posts}, \texttt{community} et \texttt{notifications}. Les captures des interfaces finales pourront compléter cette section sans remplacer les diagrammes de conception.

\section{Conclusion de la Release 1}
\label{sec:r1_conclusion}

Cette release fournit un incrément fonctionnel cohérent : accès sécurisé, information locale de tri et participation communautaire. Les diagrammes présentés décrivent les responsabilités réellement observées dans le dépôt, les échanges entre Flutter et FastAPI ainsi que les entités persistées. Cette base prépare la Release 2 consacrée à la sensibilisation, aux récompenses et à l'intégration de la poubelle intelligente.
